\section{Background and Related Work}

\paragraph{Cache Management} 
To fill the speed gap between the processor
and memory, multiple levels of caches are employed in today's computer
systems.  Cache allocation and replacement policies are important
design considerations~\cite{Smith:CSUR-1982}. 
They determine where to put a block in the cache and 
which block to be replaced when a cache miss happens,
and affect the memory system performance and enegy efficiency
significantly. A lot of research efforts have been done in this
area especially for multi-core and multi-threaded
processors where multiple applications run concurrently and compete for
the limited cache space.  For instance, various partitioning techniques
on shared caches have been proposed. 
The partitioning might be way-based, set-based, or 
block-based and be guided by the estimated impact on cache 
misses, memory latency, IPC, bandwidth or other performance
factors~\cite{Albonesi:MICRO-1999, Suh:HPCA-2002,
Kim:PACT-2004, Suh:JSC-2004, Chandra:HPCA-2005, 
Chang:ISCA-2006, Varadarajan:MICRO-2006, Hsu:PACT-2006,
Qureshi:MICRO-2006, Lin:HPCA-2008, Awasthi:HPCA-2009, 
Xie:ISCA-2009, Kaseridis:ICPP-2009, Kaseridis:HPCA-2010,
Sanchez:ISCA-2011, Manikantan:ISCA-2012,
Cook:ISCA-2013, Khan:HPCA-2014, Ye:PACT-2014,
Brock:ICPP-2015, Gupta:ICPP-2015, Subramanian:MICRO-2015,
Selfa:PACT-2017, Mittal:CSUR-2017, ElSayed:HPCA-2018, Xiao:ICPP-2019,
Chung:PACT-2019, Sun:IPDPS-2019, Pons:TPDS-2020, Park:TC-2020,
Kwon:MICRO-2023, Lee:MICRO-2025}.
They usually work with the cache placement and replacement policies to optimize
the performance, energy efficiency, fairness or QoS~\cite{Stone:TC-1992,
Sherwood:ICS-1999, Takagi:ICS-2004,
Varadarajan:MICRO-2006, Keramidas:ICCD-2007, Qureshi:ISCA-2007,
Jaleel:PACT-2008, Kharbutli:TC-2008, Liu:MICRO-2008,
Jaleel:ISCA-2010, Sanchez:MICRO-2010, 
Wu:TACO-2011, Wu:MICRO-2011, Duong:MICRO-2012, Wong:Online-2013, 
Tian:JMM-2014, Wang:MICRO-2014, Das:TACO-2015, Teran:MICRO-2016,
Beckmann:HPCA-2017, Jimenez:MICRO-2017, Kim:ASPLOS-2017,
Shi:MICRO-2019, Zhang:ICPADS-2020, Sethumurugan:HPCA-2021,
Championship:2010, Championship:2017}.

combine prefetching and replacement~\cite{Kim:ASPLOS-2017}

\paragraph{Prefetching} Prefetching is an effective technique that
can hide the long memory access latency.  It predicts future cache
misses based on the current memory access pattern (usually current
and past cache misses), and issues those predicted future misses
in advance in hope that when future misses happen, the required
data have already been fetched from lower level(s) of memory
hierarchy~\cite{Vanderwiel:CSUR-2000, Mittal:CSUR-2016}.  Prefetching
can be either software-based or hardware-based.  The software-based
approaches insert explicit prefetch instructions into the code by
programmers or compilers and usually only work for applications
with regular access patterns.  Most performance-oriented processors
today implement hardware prefetchers that can predict future cache
misses during run-time.  They usually utilize metadata within the
processor, such as PC, request's tag, delta between successive
requests, and signature of recent requests, to predict future cache
misses~\cite{Fu:MICRO-1992, Chen:TC-1995, Joseph:ISCA-1997,
Kumar:ISCA-1998, Reinman:MICRO-1999, 
Hu:HPCA-2003, Nesbit:HPCA-2004, Somogyi:ISCA-2006,
Manikantan:ICS-2008, Ishii:ICS-2009, 
Ebrahimi:HPCA-2009, Grannaes:HPEAC-2010, Wenisch:HPCA-2009,
Albericio:TACO-2012, Jain:MICRO-2013,
Shevgoor:MICRO-2015, Yu:MICRO-2015, 
Kim:MICRO-2016, Michaud:HPCA-2016,
Zeng:MEMSYS-2017, Bakhshalipour:HPCA-2018,
Hashemi:ICML-2018, Heirman:PACT-2018, Kondguli:ISCA-2018,
Bakhshalipour:HPCA-2019, Bera:MICRO-2019, Wu:MICRO-2019, Wu:ISCA-2019,
Ayers:ASPLOS-2020, Pakalapati:ISCA-2020, 
Bera:MICRO-2021, Vavouliotis:ISCA-2021, Shi:ASPLOS-2021,
Jiang:MICRO-2022, Lurbe:TC-2022,
Litz:ASPLOS-2022, Navarro:MICRO-2022, Vavouliotis:MICRO-2022,
Zhang:SC-2022,
Gerogiannis:MICRO-2023, Panda:MICRO-2023, Anisworth:ISCA-2024, 
Lin:TCAD-2024, Oh:ISCA-2024, Nath:MICRO-2024,
Ai:MICRO-2025, He:MICRO-2025, Li:ISCA-2025, 
Navarro:TC-2025, Xue:TACO-2025, Zhang:TPDS-2025, Chen:HPCA-2025,
Huang:SPAA-2026, Huang:TACO-2026, Jiang:TACO-2026,
Jamet:ISCA-2026, Lin:HPCA-2026, Wang:ISCA-2026, Xu:ASPLOS-2026,
Ye:ISCA-2026, Shao:CAL-2026,
Championship:2009, Championship:2015, Championship:2017,
Championship:2026}.

Prefetch filtering~\cite{Lin:ICCD-2001, Zhuang:ICPP-2003, Zhuang:TC-2007,
Panda:CAL-2016, Bhatia:ISCA-2019, Jamet:HPCA-2024,
Vavouliotis:HPCA-2025}.

Prefetch (aggressiveness) controller~\cite{Srinath:HPCA-2007, 
Ebrahimi:MICRO-2009, Pugsley:HPCA-2014,
Panda:TACO-2015, 
Panda:TC-2016, Sridharan:TACO-2017, Yang:TPDS-2025,
Min:CAL-2026}.

Prefetch + off-chip prediction~\cite{Bera:MICRO-2022, Bera:HPCA-2026}.

Prefetcher selection~\cite{Li:HPCA-2025}.

Bank-level parallelism aware prefetch issuing~\cite{Lee:MICRO-2009}.

Prefetch-aware resource/cache management~\cite{Ebrahimi:ISCA-2011,
Wu:MICRO-2011-2, Jain:ISCA-2018}.

Neural model of data prefetching~\cite{Shi:ASPLOS-2021} (note:already cited above).

\paragraph{Memory Access Scheduling}

Prefetch-aware scheduling~\cite{Lee:MICRO-2008}

\paragraph{DNN}

Reinforcement learning on architecture~\cite{Gerogiannis:MICRO-2023}
RL framework for prefetching~\cite{Zhang:SC-2022}
Deep learning multi-program prefetch configuration for the IBM POWER 8~\cite{Lurbe:TC-2022}
